%% ================================================================
%% SFST v30 — Neuer Supplement-Anhang: Holographische Sättigung
%% Titel: Appendix P: Holographic saturation condition and the
%%         c_1/c_2 normalization problem
%%
%% Einzufügen in sfst_supplement_v30.tex nach Appendix O (ca. Zeile 945).
%% Tier-Stufen: Tier 1 für c_1-Berechnung; Tier 2 für Proposition H;
%% Tier 3 für G_5-Ableitung aus erster Prinzipien.
%% ================================================================

\section{Appendix P: Holographic saturation condition and the
  $c_1/c_2$ normalization --- explicit assumptions and their scope}
\label{app:holographic}

This appendix unpacks the holographic saturation argument for the
distinguished radius $R_0=\tfrac{1}{2}$ in full calculational detail.
Its purpose is to make every physical identification and normalization
choice explicit, so that a reviewer can locate exactly which step
remains an assumption rather than a derived fact.  The section is
structured in three layers of decreasing rigor, in keeping with the
Tier~1/2/3 discipline of the main text.

%% ----------------------------------------------------------------
\subsection{Layer 1 (Tier~1): $c_1$ is fully determined by the
  SFST operator}
\label{ssec:c1}

The Weyl spectral counting function for the SFST operator
$\mathcal{D}$ acting on the bundle
$\mathcal{E}=V_C\otimes S(T^5_R)$ of rank $r=24$ counts eigenvalues
below a spectral cutoff~$\Lambda$:
\begin{equation}
  \label{eq:Weyl}
  N(\Lambda,R) \;=\;
  \underbrace{r}_{\text{rank}} \cdot
  \underbrace{(2\pi R)^5}_{\mathrm{Vol}(T^5_R)} \cdot
  \frac{\Lambda^5}{(2\pi)^5} \cdot
  \underbrace{\omega_5}_{\text{unit-ball vol}}
  \;=\;
  c_1 \cdot R^5 \cdot \Lambda^5,
\end{equation}
where $\omega_5 = \tfrac{8\pi^2}{15}$ is the volume of the
five-dimensional unit ball.  Inserting $r=24$:
\begin{equation}
  \label{eq:c1value}
  \boxed{c_1 = r\,\omega_5 = 24\cdot\frac{8\pi^2}{15}
              = \frac{192\pi^2}{15}
              \approx 126.33.}
\end{equation}
This value depends only on the operator definition (rank of
$\mathcal{E}$, dimensionality of $T^5$).  No free parameter enters.

\paragraph{Numerical check.}
$192/15 = 12.8$, so $c_1 = 12.8\,\pi^2 \approx 12.8 \times 9.8696
\approx 126.33$.  A reviewer can verify \cref{eq:c1value} by direct
computation.

%% ----------------------------------------------------------------
\subsection{Layer 2 (Tier~2): The saturation condition and the
  residual assumption on $G_5\cdot\Lambda^5$}
\label{ssec:saturation}

The holographic saturation argument equates the spectral entropy
$S_{\mathrm{sp}}$ (Weyl mode count) with the Bekenstein--Hawking
entropy $S_{\mathrm{BH}}$ of the torus.

\paragraph{Spectral entropy.}
Following \cref{eq:Weyl}:
\begin{equation}
  S_{\mathrm{sp}}(R,\Lambda) = c_1\,R^5\,\Lambda^5.
\end{equation}

\paragraph{Bekenstein--Hawking entropy on $T^5_R$.}
The natural codimension-one boundary of $T^5_R$ is a $T^4_R$ slice of
area
\begin{equation}
  A_4(R) = (2\pi R)^4.
\end{equation}
The Bekenstein--Hawking entropy in $d+1=6$ bulk dimensions with
Newton constant $G_5$ is
\begin{equation}
  \label{eq:SBH}
  S_{\mathrm{BH}}(R) = \frac{A_4(R)}{4\,G_5}
                     = \frac{(2\pi R)^4}{4\,G_5}.
\end{equation}

\paragraph{Saturation condition.}
Setting $S_{\mathrm{sp}} = S_{\mathrm{BH}}$ and solving for $R$:
\begin{align}
  c_1\,R^5\,\Lambda^5 &= \frac{(2\pi R)^4}{4\,G_5}\nonumber\\
  c_1\,R\,\Lambda^5   &= \frac{(2\pi)^4}{4\,G_5}\nonumber\\
  \label{eq:Rvac}
  R_{\mathrm{vac}} &= \frac{(2\pi)^4}{4\,G_5\,c_1\,\Lambda^5}.
\end{align}

\paragraph{Condition for $R_{\mathrm{vac}}=\tfrac{1}{2}$.}
Inserting $c_1=192\pi^2/15$ into \cref{eq:Rvac} and requiring
$R_{\mathrm{vac}}=\tfrac{1}{2}$:
\begin{equation}
  \label{eq:G5Lambda}
  \frac{1}{2} = \frac{(2\pi)^4}{4\,G_5 \cdot \tfrac{192\pi^2}{15}
                 \cdot \Lambda^5}
  \qquad\Longrightarrow\qquad
  G_5\,\Lambda^5 = \frac{(2\pi)^4\cdot 15}{4\cdot 192\,\pi^2\cdot 2}
               = \frac{240\,\pi^4}{1536\,\pi^2}
               = \frac{5\pi^2}{32}.
\end{equation}

\begin{tcolorbox}[colback=gray!8,colframe=black!40,title={\small
  Explicit residual assumption (Assumption~H)}]
\textbf{Assumption~H (holographic unit convention).}
The product of the five-dimensional Newton constant and the fifth power
of the spectral cutoff satisfies
\[
  G_5\,\Lambda^5 \;=\; \frac{5\pi^2}{32} \;\approx\; 1.5421.
\]
Under this assumption, the saturation condition
$S_{\mathrm{sp}}=S_{\mathrm{BH}}$ uniquely determines
$R_{\mathrm{vac}}=\tfrac{1}{2}$.
\end{tcolorbox}

Assumption~H is a Tier~2 statement.  It is internally consistent and
physically motivated, but it is not derived from a deeper SFST
principle in the current manuscript.  Its precise content is the
dimensionless product $G_5\Lambda^5$, not the individual values of
$G_5$ or $\Lambda$.

%% ----------------------------------------------------------------
\subsection{Layer 3 (Tier~3): Why $G_5\,\Lambda^5 = 5\pi^2/32$
  is not yet derivable}
\label{ssec:G5open}

Assumption~H involves two physically distinct objects:
$G_5$ (a gravitational coupling) and $\Lambda$ (a spectral cutoff
scale of the SFST operator).  Three routes toward deriving their
product are explored here; all remain incomplete.

\paragraph{Route A: Kaluza--Klein reduction from $G_4$.}
The standard KK relation between Newton constants in adjacent
dimensions is
\[
  G_4 = \frac{G_5}{2\pi R_{\mathrm{vac}}}.
\]
Inserting $R_{\mathrm{vac}}=\tfrac{1}{2}$:
$G_4 = G_5/\pi$, hence $G_5 = \pi\,G_4$.
In four-dimensional Planck units with $G_4=1$:
$G_5 = \pi$.  This gives
\[
  G_5\,\Lambda^5 = \pi\,\Lambda^5,
\]
which reproduces Assumption~H only if $\Lambda^5=5\pi/32$, i.e.\
$\Lambda=(5\pi/32)^{1/5}\approx 0.859$.  This is a specific numerical
prediction for the spectral cutoff in 4D Planck units, but it is not
independently motivated within SFST at present.

\paragraph{Route B: Direct 5D Planck units.}
Setting $G_5=1$ (5D Planck units, $\ell_{\mathrm{P5}}=1$):
\[
  G_5\,\Lambda^5 = \Lambda^5.
\]
Assumption~H then becomes $\Lambda=(5\pi^2/32)^{1/5}\approx 1.090$.
Again a specific cutoff prediction, not derived from SFST dynamics.

\paragraph{Route C: SFST natural units.}
If one defines SFST natural units by the condition
$R_{\mathrm{vac}}=\tfrac{1}{2}$ as a normalization convention, then
Assumption~H is a tautology: the units are chosen so that
$G_5\Lambda^5=5\pi^2/32$.  This is a permissible organizational
choice, but it offers no physical content and does not explain why
$R_{\mathrm{vac}}=\tfrac{1}{2}$ is preferred over other values.

\paragraph{Summary of Layer~3.}
None of Routes A--C derives $G_5\Lambda^5=5\pi^2/32$ from an
independent SFST principle.  Route~A is the physically most natural
interpretation; it converts the problem into a prediction for the
spectral cutoff $\Lambda$ in Planck units, which could in principle be
tested against an independent determination of the SFST energy scale.
Routes B and C are self-consistent but do not add physical content.

%% ----------------------------------------------------------------
\subsection{Relation to the two earlier conditional assumptions
  A1 and A2}
\label{ssec:relAtoH}

Previous versions of the manuscript grounded $R_0=\tfrac{1}{2}$ in
two separate assumptions: T-duality (A1) and $\alpha'=\tfrac{1}{4}$
(A2).  Assumption~H replaces those two by a single physical condition:
holographic saturation of the SFST vacuum.

The replacement has clear advantages:
\begin{itemize}[leftmargin=1.5em]
  \item Winding modes are not required (A1 was a structural postulate).
  \item The factor $1/4$ in the Bekenstein--Hawking formula is not
        postulated separately; it is the standard gravitational
        coefficient.
  \item The condition is falsifiable in principle: if a future
        determination of the SFST energy scale sets
        $G_5\Lambda^5\neq 5\pi^2/32$, Assumption~H is ruled out.
\end{itemize}

The replacement also has a precise limitation:
\begin{itemize}[leftmargin=1.5em]
  \item The single remaining free input is the dimensionless product
        $G_5\Lambda^5$.  Until this is fixed by an independent
        dynamical or gravitational argument, Assumption~H is a
        Tier~2 conditional statement, not a Tier~1 theorem.
  \item The $5$D Newton constant $G_5$ is dimensionally distinct from
        the $4$D Newton constant $G_4$ that enters gravitational
        phenomenology.  The precise connection requires a
        compactification argument (Route~A above) that is not yet
        completed within SFST.
\end{itemize}

%% ----------------------------------------------------------------
\subsection{What the manuscript must display transparently}
\label{ssec:transparency}

In keeping with the evidentiary discipline of this supplement, the
following items must remain explicitly visible in the main text and
this appendix:

\begin{enumerate}[label=(\roman*)]
  \item $c_1 = 192\pi^2/15$ is fully determined by the SFST operator
        (Tier~1).
  \item The identification $c_2 = (2\pi)^4/(4G_5)$ uses the
        Bekenstein--Hawking formula with a specific area choice
        $A_4=(2\pi R)^4$ (physical identification, not a theorem).
  \item Assumption~H, namely $G_5\Lambda^5=5\pi^2/32$, is the unique
        residual normalization condition required to obtain
        $R_{\mathrm{vac}}=\tfrac{1}{2}$ from the saturation equation
        (Tier~2 conditional).
  \item The derivation of $G_5\Lambda^5$ from first principles is an
        open problem; the three routes explored above reduce it to a
        prediction for the spectral cutoff or to a unit choice,
        neither of which is currently derivable from SFST dynamics
        alone (Tier~3 conjecture).
\end{enumerate}

This four-point transparency checklist is the minimal record a referee
would need to locate the logical boundary between what is calculated,
what is assumed, and what remains open in the holographic radius
argument.
